\chapter{Methodology}
\label{ch:methodology}

\section{Development Methodology}

Echo was built using \textbf{iterative prototyping}. This means each part
of the system was built on its own, tested on its own, and only then
connected to the rest, instead of designing the whole system on paper
first and building it all in one go. This approach was chosen on purpose,
because Echo's design is a strict chain: sound goes in, the server checks
it, and the laptop gets logged in. A problem in any one part would block
everything after it. Checking each part before building on top of it kept
problems small and easy to find. The seven build stages, in the order they
happened, were:

\begin{enumerate}[label=\textbf{Stage \arabic*.}, leftmargin=2.2cm]
  \item \textbf{Does the sound idea even work?} Before writing any server
  code, a simple test page (\texttt{tests/\allowbreak ultrasonic-\allowbreak auth-\allowbreak test.html}) was
  built to answer one question: can a real laptop speaker and a real phone
  microphone actually send a signal at 18 to 20\,kHz (too high-pitched for
  most people to hear), and at what distance does it stop working? This
  test also produced the first written list of what the server and phone
  needed to do.
  \item \textbf{Build the server first.} The backend was built before any
  screen the user would see, because the screens needed something to talk
  to. The full database was set up (users, phones, login attempts,
  recovery codes, login history), and every part of the API was written
  and tested as it was added, along with the live-update connection.
  \item \textbf{Build the sound layer and the phone app.} With the server
  ready, the laptop login page was built to connect to \texttt{ggwave} (the
  sound library) and send the one-time code. Then the phone app was built.
  It listens through the microphone, decodes the sound, and shows an
  approve or deny screen. Both sides were tested together in the same room
  first, then at increasing distances. A quiet background check was added
  so the phone can confirm, before showing anything to the user, that a
  decoded code actually belongs to the right account.
  \item \textbf{Add the cryptography.} Once the sound channel worked
  reliably, the sign-up process was built so the phone creates a private
  signing key and registers the matching public key with the server. The
  phone's signing code was written, and the server's signature-checking
  code was finished. The whole login process was then tested from start to
  finish: type username, hear sound, tap approve, get logged in. The
  automated tests were expanded to cover attack scenarios directly.
  \item \textbf{Add extra security and polish.} With the main flow
  working, the remaining features were added: a rate-limited way to
  recover an account, the ability to remove a registered phone, a one-time
  claim step so only the right laptop gets logged in, a live status
  display on the login page (transmitting, then phone heard it, then
  approved), and a normal audible backup for hardware that cannot play
  very high-pitched sound.
  \item \textbf{Test how well it actually works.} A test tool was built to
  run controlled trials, recording whether the phone picked up the code,
  whether the login finished, how long it took, the distance, and how
  noisy the room was, across a range of distances and noise levels
  (Section~\ref{sec:evaluation}).
  \item \textbf{Write it up.} The final stage produced the API reference,
  the description of how the system is built, the test results, and the
  security discussion presented in this report.
\end{enumerate}

While building the system, the account-recovery method originally planned
(numeric one-time recovery codes) was replaced with a rate-limited email
link instead, and the sign-up process was extended with a QR-code scanner
on the phone to make it faster. Both of these changes are normal for
iterative prototyping: a working system shows up usability problems that a
plan on paper does not. Both changes are discussed further in
Section~\ref{sec:design-evolution}.

\section{System Architecture}
\label{sec:architecture}

Echo has three parts that work together: a \emph{laptop page} (the device
being logged into, running in a normal web browser), a \emph{phone app}
(the trusted, registered device, installed from a browser), and a
\emph{server}, built with Node.js, that sits between them and stores all
the data. Figure~\ref{fig:architecture} shows the full login process, step
by step.

\begin{figure}[htbp]
\centering
\resizebox{\textwidth}{!}{%
\begin{tikzpicture}[
  node distance=1.4cm and 2.6cm,
  box/.style={draw, rounded corners, minimum width=3.4cm, minimum height=1cm, align=center, font=\small},
  arr/.style={-{Latex[length=2.5mm]}, thick},
  lbl/.style={font=\scriptsize, align=center}
]
\node[box] (laptop) {Laptop browser\\(login page)};
\node[box, right=of laptop, xshift=2.0cm] (server) {Echo server\\Node.js + Express + SQLite};
\node[box, right=of server, xshift=2.0cm] (phone) {Phone app\\(listens + signs)};

\draw[arr] (laptop.north) -- ++(0,0.8) -| node[lbl, above, pos=0.25, yshift=1pt] {1. Ask for a one-time code} (server.north);
\draw[arr] (server.north) ++(0,0.8) -| node[lbl, above, pos=0.75, yshift=1pt] {2. Sends back the code} (laptop.north);

\draw[arr, dashed] (laptop.north) to[out=60, in=120, looseness=1.6]
  node[lbl, above, yshift=2pt] {3. Code played as sound, 18 to 20\,kHz\\about 1 to 2\,m away} (phone.north);

\draw[arr] (phone.south) -- ++(0,-0.8) -| node[lbl, below, pos=0.75, yshift=-1pt] {4. Phone sends back a signed approval} (server.south);

\draw[arr] (server.south) ++(0,-0.8) -| node[lbl, below, pos=0.25, yshift=-1pt] {5. Server confirms: logged in} (laptop.south);
\end{tikzpicture}%
}
\caption{The Echo login process: the server creates a one-time code, the laptop plays it as sound, the phone signs it once the user approves, and the server tells the laptop it is logged in.}
\label{fig:architecture}
\end{figure}

The technology was chosen to keep things simple and quick to build, not to
handle millions of users. Node.js with Express handles the web requests;
SQLite, using a feature already built into Node.js, stores users, phones,
login attempts, recovery codes, and login history, without needing a
separate database program running; a small library called \texttt{ws}
handles the live WebSocket connection; both the laptop and phone pages are
plain JavaScript with no front-end framework, since neither one needs
enough on-screen state management to justify one; and \texttt{ggwave},
which runs in the browser, is what turns data into sound and back again on
both ends. Table~\ref{tab:stack} lists the technology used at each layer.

\begin{table}[htbp]
\centering
\caption{Technology used at each layer}
\label{tab:stack}
\begin{tabular}{|L{4cm}|L{9cm}|}
\hline
\textbf{Layer} & \textbf{Technology} \\
\hline
Server & Node.js version 22.5 or later, with Express 4 \\
\hline
Database & SQLite, built into Node.js (WAL mode for reliability) \\
\hline
Live updates & \texttt{ws} (WebSocket) \\
\hline
Sound sending/receiving & \texttt{ggwave}, running in the browser \\
\hline
Cryptography & WebCrypto (on the phone); Node's built-in crypto tools (on the server) \\
\hline
Phone key storage & IndexedDB, a browser storage area the key cannot be copied out of \\
\hline
Phone install & Web App Manifest + Service Worker (lets it work offline) \\
\hline
Web pages & Plain JavaScript, HTML, CSS, no framework \\
\hline
Hosting & Render.com, set up through a \texttt{render.yaml} file \\
\hline
\end{tabular}
\end{table}

\section{How the Login Process Works}
\label{sec:methodology-crypto}

\subsection{Signing Up}

Signing up only happens once. The user picks a username
(\texttt{POST /api/signup}) and gets back a one-time sign-up code. The
phone, once it has that code (at first typed in by hand, later scanned as
a QR code, see Section~\ref{sec:design-evolution}), creates a private
signing key and a matching public key, right there in the browser, and
sends the public key to the server along with the sign-up code
(\texttt{POST /api/enroll}). The server checks that the key is the right
type, marks the sign-up code as used, and saves the public key against the
user's account. The private key is created in a way that cannot be copied
out, so it only ever exists inside the phone's browser and can never be
read, exported, or leaked out by any code, including Echo's own.

\subsection{Logging In}

\begin{enumerate}[label=\arabic*.]
  \item The user types their username on the laptop. The server creates a
  random one-time code (a "nonce") tied to that one login attempt
  (\texttt{POST /api/login/start}), good for 30 seconds and usable only
  once.
  \item The laptop turns the code into sound using \texttt{ggwave} and
  plays it at 18 to 20\,kHz. If nothing happens, it automatically tries
  again, up to three times.
  \item The phone, always listening, picks up the code and quietly checks
  with the server that this code belongs to an account it is actually
  registered to, before showing anything on screen. This stops a phone
  from ever asking its owner to approve someone else's login.
  \item If that check passes, the phone shows an approve or deny screen.
  If the user approves, the phone signs a short message containing the
  code and its own device ID, using its private key, and sends that
  signature to the server (\texttt{POST /api/login/verify}).
  \item The server looks up the phone's saved public key, checks that the
  signature matches, checks the code has not expired or been used before,
  and if everything checks out, marks the code as used and sends a
  message over the live connection telling the laptop it is approved,
  along with a one-time claim token.
  \item The laptop swaps that claim token for a proper login cookie
  (\texttt{POST /api/\allowbreak session/\allowbreak claim}), which is set with safety flags
  (\texttt{HttpOnly}, \texttt{SameSite=Strict}, and \texttt{Secure} in
  production). The claim token can only be used once, so only the laptop
  that started this exact login can use it.
\end{enumerate}

The exact message the phone signs, \texttt{echo-v1|<code>|<deviceId>},
matters more than it looks. Starting it with a version label means the
system could change later without confusion, and including the phone's own
device ID, not just the one-time code, stops a signature made for one
phone from being reused against a different phone's record, even if an
attacker somehow got hold of it. Table~\ref{tab:nonce} lists the rules the
one-time code follows, and every one of these rules is directly checked by
the automated tests in Section~\ref{sec:testing}.

\begin{table}[htbp]
\centering
\caption{Rules for the one-time login code}
\label{tab:nonce}
\begin{tabular}{|L{4.5cm}|L{8.5cm}|}
\hline
\textbf{Property} & \textbf{Rule} \\
\hline
Randomness & A random 96-bit value, picked by the server \\
\hline
Time limit & Valid for 30 seconds \\
\hline
Number of uses & Works once only; marked as used right after \\
\hline
Scope & Tied to one login attempt and one username \\
\hline
Signed message & \texttt{echo-v1|<code>|<deviceId>}, ties the code to one specific phone \\
\hline
How it travels & Not secret: sent as sound (laptop to phone), then again as plain text when the phone checks in. This is safe because the code can only be used once \\
\hline
\end{tabular}
\end{table}

\section{Why This Design Is Secure}

Three design choices do most of the security work here, and each one maps
to a real decision in the code. First, \emph{the one-time code is public
but can only be used once}, so recording and replaying the login sound
gets an attacker nothing: the first successful check uses up the code, and
the server rejects any attempt to use it again, which is tested directly
in Section~\ref{sec:testing}. Second, \emph{the private key never leaves
the phone}: because it is created in a way that cannot be exported, there
is no way, not even through a bug in Echo's own code, to read the raw key
back out of the browser. Third, \emph{the signature always goes straight
from the phone to the real Echo server, never through the sound itself or
through whatever page the laptop happens to be showing}. This is what
stops a fake copy of the login page from working: even if a victim's
laptop is pointed at a fake site playing the same sound, the victim's
phone still only sends its signature to the real Echo server, because
that destination is decided by the phone's own code, not by the sound it
heard. A fake site can record and replay the sound, but it cannot make the
victim's phone send its signature anywhere else. Distance adds a fourth,
weaker layer of protection: because an 18 to 20\,kHz sound fades quickly in
air and does not travel well through walls, an attacker outside the room
cannot trigger the sound part of the process at all. Chapter~\ref{ch:conclusion}
explains why this is a helpful sign of nearby presence rather than a
cryptographic guarantee, and how that is different from a proper defence
against a relay attack.

\section{How the Testing Was Planned}
\label{sec:methodology-eval}

Two separate testing tools were built, because there are two different
things that need to work: the cryptography needs to be correct, and the
sound channel needs to be reliable in the real world. Correctness was
checked with a 31-check automated test suite (Section~\ref{sec:testing})
that talks to the server directly over the network, using real signing
keys but no actual audio hardware, so it runs in seconds and can be
repeated after every change without needing a quiet room. Sound
reliability was checked separately, by hand, with controlled trials
(Section~\ref{sec:evaluation}) at three distances (0.3\,m, 1\,m, 2\,m) and
in three noise conditions (quiet, speech, music), with at least 30 trials
for each combination, recording whether the phone picked up the code,
whether the login finished, and how long it took. This matches the target
set out in the objectives (Section~\ref{ch:introduction}): at least 90\%
success at 1\,m in a quiet room.
